\subsection{Technical terminology}\label{subsec:techical-research}
This subsection provides context on some of the terms this thesis uses.

\subsubsection*{CMS}\label{subsubsec:cms}
Content Management System is type of software that provides simple interface to create, manage, and publish digital
content without using a programming language~\parencite{oracle_cms}.
The software must provide sufficient tools for management of large datasets from different sources such as Social Media,
Internet of Things, and other~\parencite{kandepu_cms}.

\subsubsection*{AI}\label{subsubsec:ai}
In this thesis the terms LLM and AI are used interchangeably and both refer to software that returns text response from natural language request.
Large Language Models are transformer-based language models that are pre-trained on web-scale corpora.
These models are capable of obtaining world knowledge, generating code, following instructions, self-improve, and other~\parencite{minaee2025_llm_suvey}.
The proposed solution works with LLMs to generate structured text responses~\parencite{cloudflare_llm}.
LLMs can be simplified to a function that takes text input written in natural language and returns structured text response,
which can include parts written in natural language.

\subsubsection*{Framework}\label{subsubsec:framework}
Represents a large code base on which multiple developers collaborate~\parencite[6]{rahman2011}.
The code is structured for maintainability, reducing repetition, and onboarding new developers.
Frameworks give strong starting point for problems in target domain and handle common functionality for the domain.

\subsubsection*{Library}\label{subsubsec:library}
Libraries are collection of reusable software components for given domain.
The software components defined in library may be classes or procedures~\parencite[sec.~2.2, p.~8]{wegner1990}.
Libraries are not used until their functionality is explicitly called.
On the other hand, frameworks provide starting point and directly call user code.

\subsubsection*{Full-stack Web Development}\label{subsubsec:fullstack-web-development}
The Web development can be simplified into two distinct parts, back-end and front-end code.
Back-end code implements business logic on the server, which includes, but it is not limited to request handling,
file serving, database communication, and server side rendering.
Front-end code is served to the user as HTML, CSS, and JavaScript files that are then parsed and executed in the browser.
Full-stack developer works on both front-end interfaces and back-end procedures~\parencite[Full-stack Web Developer]{w3schools_fullstack}.